\chapter{Database Design}

\section{Entity Model}

The data model is designed around durable business evidence. It must support
merchant readiness, transaction state, refund state, notification delivery,
external outcome evidence, reconciliation, audit history, and human account
access without losing traceability.

\begin{landscape}
\begin{figure}[H]
\centering
\begin{tikzpicture}[node distance=11mm and 14mm]
  \node[box, text width=32mm] (merchant) {Merchant};
  \node[box, above left=of merchant, text width=32mm] (internalUser) {Internal user};
  \node[box, below left=of merchant, text width=34mm] (onboarding) {Onboarding case};
  \node[box, below=of onboarding, text width=34mm] (credential) {Merchant credential};
  \node[box, right=of merchant, text width=34mm] (payment) {Payment transaction};
  \node[box, above right=of payment, text width=34mm] (merchantUser) {Merchant dashboard user};
  \node[box, below right=of payment, text width=34mm] (refund) {Refund transaction};
  \node[box, right=of payment, text width=34mm] (webhook) {Webhook event};
  \node[box, below=of webhook, text width=34mm] (attempt) {Webhook delivery attempt};
  \node[box, above=of webhook, text width=34mm] (callback) {Callback evidence};
  \node[box, below=of refund, text width=34mm] (recon) {Reconciliation record};
  \node[box, below=of merchant, text width=34mm, xshift=44mm] (audit) {Audit log};

  \draw[arrow] (merchant) -- (onboarding);
  \draw[arrow] (merchant) -- (credential);
  \draw[arrow] (merchant) -- (merchantUser);
  \draw[arrow] (merchant) -- (payment);
  \draw[arrow] (payment) -- (refund);
  \draw[arrow] (payment) -- (webhook);
  \draw[arrow] (webhook) -- (attempt);
  \draw[arrow] (callback) -- (payment);
  \draw[arrow] (callback) -- (refund);
  \draw[arrow] (callback) -- (recon);
  \draw[arrow] (payment) -- (recon);
  \draw[arrow] (refund) -- (recon);
  \draw[arrow] (internalUser) -- (onboarding);
  \draw[arrow] (internalUser) -- (audit);
  \draw[arrow] (merchantUser) -- (audit);
  \draw[arrow] (merchant) -- (audit);
\end{tikzpicture}
\caption{Logical entity relationship overview. The model centers on merchants,
transactions, delivery evidence, and operational traceability.}
\end{figure}
\end{landscape}

\section{Core Entities}

\begin{longtable}{L{0.26\textwidth}L{0.64\textwidth}}
\caption{Core business entities}\\
\toprule
\textbf{Entity} & \textbf{Purpose} \\
\midrule
\endfirsthead
\toprule
\textbf{Entity} & \textbf{Purpose} \\
\midrule
\endhead
Merchant & Stores the merchant identity, operating profile, contact data, notification destination, settlement information, and current operating status. \\
Onboarding Case & Stores merchant-submitted readiness information, review findings, approval or rejection outcomes, and review history. \\
Merchant Credential & Stores merchant integration credential metadata and protected secret material used for signed transaction requests. \\
Merchant Dashboard User & Stores human merchant accounts that sign in to the merchant dashboard and are scoped to one merchant. \\
Internal User & Stores internal administrative and operational accounts used for management and review workflows. \\
Payment Transaction & Stores payment attempts, money amounts, state, expiration information, merchant references, and externally linked payment evidence. \\
Refund Transaction & Stores refund requests and their linkage to successful original payments. \\
Webhook Event & Stores merchant notification payloads, delivery state, retry metadata, and notification history. \\
Webhook Delivery Attempt & Stores individual delivery attempts so the system can preserve delivery evidence over time. \\
Callback Evidence & Stores raw and normalized external outcome evidence for payments and refunds. \\
Reconciliation Record & Stores manual-review evidence when external outcomes do not align cleanly with internal state. \\
Audit Log & Stores durable review and administrative evidence for sensitive human actions. \\
\bottomrule
\end{longtable}

\section{Data Invariants}

\begin{longtable}{L{0.34\textwidth}L{0.56\textwidth}}
\caption{Important data invariants}\\
\toprule
\textbf{Invariant} & \textbf{Reason} \\
\midrule
\endfirsthead
\toprule
\textbf{Invariant} & \textbf{Reason} \\
\midrule
\endhead
Each merchant has one stable public merchant identity. & External integrations and both dashboards need one durable merchant anchor. \\
Only one active merchant integration credential exists at a time. & Signed requests should resolve to one authoritative active credential for each merchant. \\
Each merchant dashboard user belongs to exactly one merchant. & Merchant dashboard access must remain tenant-scoped and must not leak across merchants. \\
Only one active pending payment can exist for the same merchant order at a time. & This prevents duplicate active checkout attempts for one business order. \\
Refund requests are unique per merchant refund reference. & The system needs business-level idempotency for merchant refund initiation. \\
At most one successful full refund can exist for one successful payment. & The project implements full-refund-only behavior. \\
Notification delivery history is preserved separately from notification state. & Final notification status and the evidence that led to it must both remain available. \\
Protected secrets and plaintext passwords are never returned through ordinary read views. & Human access and merchant integration credentials must remain protected even when metadata is viewable. \\
\bottomrule
\end{longtable}

\section{Data Evolution}

The data model evolves by extending business capabilities rather than by
rewriting existing evidence. As the system grows, new entities or attributes
can be introduced for additional merchant lifecycle information, expanded
notification channels, or richer reconciliation support while preserving
historical transaction and audit integrity.
